\subsection{Testing the flag-wide band and skew-aware intervals}
\label{sec:repair}
A separate fixed-design follow-up evaluates 30 and 50 cells, using 2{,}000
fresh studies per cohort/pool/shift: 32{,}000 studies and 160{,}000 interval
evaluations. It compares percentile intervals with 200 and 2{,}000 resamples,
bias-corrected and accelerated (BCa) and bootstrap-$t$ intervals with 2{,}000
resamples, and Welch. These skew-aware methods follow established bootstrap
accuracy work \citep{diciccio1996}. All methods use the same study samples;
bootstrap methods share resamples. Coverage uses the fixed-pool target.
Appendix~\ref{app:repair} specifies implementation and numerical checks.

\begin{table}[htbp]
\centering\small
\caption{Fixed-pool evaluation inside the flag-wide band and at the reporting
boundary. Each cohort/pool/count/shift uses 2{,}000 studies; all methods share
samples. Bootstrap budgets follow method names. $C_{.5}$ is coverage at the
halving, $R_1$ null rejection (one minus null coverage), and $D_{.5}$ detection
at the halving, all in percent. Nominal coverage is 95\%; the historical
alternative-coverage minimum is 90\% and detection requirement 80\%.
Maximum Monte Carlo SE is 1.12 percentage points; Wilson intervals and paired
comparisons are supplied with the aggregate results. All intervals are finite.}
\label{tab:repair}
\begin{tabular}{@{}llrrrrrr@{}}
\toprule
 & & \multicolumn{3}{c}{30 cells: flag-wide} & \multicolumn{3}{c}{50 cells: report} \\
\cmidrule(lr){3-5}\cmidrule(l){6-8}
Cohort / pool & Interval & $C_{.5}$ & $R_1$ & $D_{.5}$ & $C_{.5}$ & $R_1$ & $D_{.5}$ \\
\midrule
SMC / pooled & Percentile, 200 & 80.5 & 31.0 & 53.6 & 87.8 & 22.8 & 48.9 \\
 & Percentile, 2000 & 82.1 & 30.9 & 52.6 & 88.4 & 22.2 & 48.0 \\
 & BCa, 2000 & 77.5 & 29.5 & 47.9 & 83.7 & 21.6 & 43.9 \\
 & Bootstrap-$t$, 2000 & 79.7 & 28.1 & 48.0 & 86.3 & 20.8 & 43.2 \\
 & Welch & 84.3 & 30.5 & 53.1 & 90.1 & 22.0 & 48.3 \\
\addlinespace
SMC / reference & Percentile, 200 & 89.1 & 14.2 & 67.9 & 91.3 & 9.5 & 77.8 \\
 & Percentile, 2000 & 90.0 & 13.8 & 67.0 & 92.1 & 8.5 & 76.3 \\
 & BCa, 2000 & 89.8 & 11.8 & 57.8 & 92.0 & 7.8 & 68.9 \\
 & Bootstrap-$t$, 2000 & 91.2 & 9.0 & 53.8 & 92.6 & 6.1 & 65.9 \\
 & Welch & 90.4 & 13.7 & 68.4 & 92.5 & 8.6 & 78.3 \\
\addlinespace
KUL3 / pooled & Percentile, 200 & 78.1 & 40.5 & 56.4 & 83.4 & 31.8 & 49.7 \\
 & Percentile, 2000 & 79.1 & 40.1 & 55.8 & 84.4 & 30.9 & 49.9 \\
 & BCa, 2000 & 72.6 & 40.3 & 54.4 & 77.5 & 32.2 & 48.4 \\
 & Bootstrap-$t$, 2000 & 75.3 & 40.4 & 55.2 & 79.9 & 31.8 & 49.7 \\
 & Welch & 83.7 & 39.4 & 56.1 & 88.2 & 30.7 & 49.8 \\
\addlinespace
KUL3 / reference & Percentile, 200 & 78.8 & 26.3 & 48.5 & 86.3 & 18.7 & 47.7 \\
 & Percentile, 2000 & 79.4 & 25.7 & 47.9 & 87.3 & 18.2 & 46.6 \\
 & BCa, 2000 & 77.3 & 24.0 & 42.2 & 85.4 & 17.0 & 39.5 \\
 & Bootstrap-$t$, 2000 & 78.4 & 22.9 & 43.0 & 86.2 & 16.1 & 40.6 \\
 & Welch & 80.5 & 26.3 & 49.7 & 88.8 & 18.6 & 48.9 \\
\addlinespace
\bottomrule
\end{tabular}
\end{table}


At 30 cells, the original interval has 78.1--89.1\% alternative coverage,
14.2--40.5\% null rejection and 48.5--67.9\% detection. Flagging its width
does not establish calibration. This evaluates one interior point, not the
whole 20--49-cell band. At 50 cells, increasing the percentile budget reduces
null rejection by 0.5--1.0 percentage points (paired Monte Carlo SEs
0.28--0.34 points), much less than its excess over nominal.

The strongest null improvement occurs for bootstrap-$t$ in SMC/reference:
rejection falls from 8.45\% for percentile with 2{,}000 resamples to
6.05\% [5.09, 7.18], a paired reduction of 2.40 points (SE 0.40).
Alternative coverage is 92.6\%, but detection falls from 76.3\% to 65.9\%.
The improvement does not transfer uniformly: in KUL3/pooled, BCa and
bootstrap-$t$ have 32.2\% and 31.8\% null rejection, against 30.9\% for
the equal-budget percentile interval. All methods return intervals, so these
comparisons are not driven by abstention. No tested method meets the historical
coverage/detection requirements together with nominal null calibration at
either count. This bounds the tested repair attempt without selecting a new rule.


\paragraph{Changing the resampling unit.}
A separate comparison with two held-out patients also tests cluster and
paired patient-$t$ intervals (Appendix~\ref{app:substitution}). At 50 cells,
patient-$t$ gives 2.77--6.29\% null rejection but only 4.21--8.02\%
detection among returned intervals; alternative abstention is 0.35--34.55\%.
It equally weights shared patients, changing the cell-weighted target.
Its saved coverage uses the sampled-reference benchmark, so it cannot
establish a fixed-pool repair. Full rates and abstentions for both clustered
methods are supplied; these results concern this design, not patient-level
inference generally.

\paragraph{Source shape and arm balance.}
A separate sensitivity study uses 2{,}000 studies and 2{,}000 percentile
resamples per setting (Appendix~\ref{app:substitution}). At 50 tumor cells,
empirical-pool null rejection is 9.70--31.80\%; removing zeros and rescaling
the positive counts to the original mean gives 7.10--12.65\%, while a
mean/variance-matched Gaussian gives 5.45\%. Keeping the empirical pools but
changing the reference count from 800 to 50 gives 7.45--11.45\%.
This implicates source shape and arm imbalance without isolating their
individual contributions. Under the null, equal-sized independent arms
have an exactly symmetric difference by exchangeability, even when the
source is skewed. Balancing also changes variance and the reference sample
size; alternative detection was not evaluated for this design change, so
these rates do not justify discarding reference cells or establish a repair.
